\documentclass[]{article}

\usepackage{amsmath}
\usepackage{multirow}
\usepackage{natbib}
\usepackage{graphicx} 


\begin{document}

We present the results of our analysis, beginning with an assessment of measurement noise and its mitigation. We then investigate the mechanistic drivers of transport, including tracer trapping and velocity memory, before characterizing the overall transport regimes through mean squared displacement scaling and ergodicity metrics.

\subsection{Noise sensitivity and data filtering}

To ensure the robustness of our findings against measurement uncertainty, we first quantified the impact of the simulated 0.02 m Gaussian noise on the velocity statistics. The analysis of single-step displacement increments revealed a strong, configuration-dependent sensitivity. For sparse vortex configurations ($N=5$ and $N=10$), where the characteristic physical displacements are small, the added noise significantly contaminated the tails of the increment distribution, artificially inflating estimates of the tail index. In contrast, for denser configurations ($N=20$ and $N=40$), the signal-to-noise ratio is much higher, and the noise has a progressively smaller effect on the tail statistics.

Based on these findings, we implemented a targeted data processing strategy. For the dense configurations ($N=20$ and $N=40$), a Savitzky-Golay filter was applied to the noisy position data. A sensitivity analysis showed that a filter window of 7 time steps provided an optimal balance, effectively suppressing high-frequency noise without distorting the underlying heavy-tailed dynamics. For the sparse configurations ($N=5$ and $N=10$), the signal-to-noise ratio was too low for effective filtering. Therefore, all subsequent mechanistic analyses for these two cases were performed exclusively on the noise-free (`true`) data to avoid confounding artifacts. This approach ensures that our conclusions are based on the physical dynamics of the system rather than measurement noise.

\subsection{Tracer trapping and residence time statistics}

To investigate the role of trapping in shaping transport, we analyzed the distribution of residence times, $P(\tau_{res})$. A tracer was considered "trapped" during time intervals when its speed fell below the 25th percentile of the speed distribution for its respective $N$ configuration. The tail of the residence time distribution was modeled as a power law, $P(\tau_{res}) \sim \tau_{res}^{-\gamma}$, with the exponent $\gamma$ estimated using maximum likelihood.

Our results show a clear trend in trapping behavior with increasing vortex density. For the sparse system $N=10$, the distribution exhibits a power-law tail with an exponent $\gamma \approx 2.1 \pm 0.3$. As the density increases, the tail becomes heavier, with $\gamma \approx 1.8 \pm 0.2$ for $N=20$ and $\gamma \approx 1.6 \pm 0.2$ for $N=40$. This indicates that while tracers in sparser systems can be trapped, the probability of extremely long trapping events increases significantly in the more chaotic, denser systems. For the nearly integrable $N=5$ case, tracers can become quasi-permanently bound to stable vortex structures, leading to a distribution dominated by a few extremely long events for which a power-law fit is ill-defined.

Comparison with a power-law model with an exponential cutoff confirmed that for the densest configurations ($N=20$ and $N=40$), a pure power-law provides a sufficient description of the data within the simulation's time window. This suggests that the heavy-tailed trapping statistics are a genuine feature of the chaotic dynamics in these systems. Physically, this corresponds to a transition from long-term trapping in stable regions at low $N$ to a regime of frequent, intermittent capture and release by transient vortex structures at high $N$.

\subsection{Velocity correlations and Hamiltonian memory}

We examined the velocity autocorrelation function (VACF) to directly probe the system's memory, comparing the point-vortex dynamics to the memoryless L\'evy walk null model. The VACF for the point-vortex system reveals the presence of significant Hamiltonian memory. For all vortex densities, the VACF exhibits a rapid initial decay followed by a persistent, positive correlation tail that extends for a significant duration before crossing zero at approximately 0.5 s. The rate of this decay is fastest for the most chaotic system ($N=40$) and slowest for the most regular one ($N=5$), reflecting the rate of velocity decorrelation in the flow.

This behavior stands in stark contrast to the L\'evy walk model. By construction, the L\'evy walk is a memoryless process, and its VACF decays to zero within 2-3 time steps and remains there. This fundamental difference highlights a key physical feature that distinguishes the deterministic Hamiltonian flow from its stochastic approximation: the point-vortex system retains a memory of its velocity over time, a direct consequence of the conservative dynamics.

These findings were corroborated by computing the mutual information between velocities at different time lags, which captures both linear and non-linear correlations. The mutual information for the point-vortex system also showed a slow decay, remaining significantly above the noise floor for much longer than in the L\'evy walk model. This confirms that the point-vortex system possesses non-Markovian velocity statistics that are entirely absent in the stochastic null model.

\subsection{Non-stationarity and ergodicity breaking}

To assess the temporal homogeneity and ergodicity of the transport process, we analyzed the time-averaged MSD (TAMSD) for individual trajectories. A sliding-window analysis of the local anomalous exponent, $\alpha(t)$, reveals that the transport is non-stationary for all configurations. For $N=5$, $\alpha(t)$ fluctuates around 2.0, consistent with persistent near-ballistic motion. For $N=20$ and $N=40$, $\alpha(t)$ is highly variable and shows a tendency to decrease over time, suggesting a crossover from an initial ballistic phase to a subsequent superdiffusive one as tracers explore the complex vortex field.

We quantified the degree of ergodicity breaking using the parameter $EB(\Delta)$, which measures the ensemble variance of the TAMSD. For the point-vortex system, the $EB$ parameter increases monotonically with vortex density, from $EB \approx 0.08$ for $N=5$ to $EB \approx 0.45$ for $N=40$. This indicates that as the system becomes more chaotic, the transport dynamics become more heterogeneous across the ensemble of tracers, leading to a greater departure from ergodicity.

Comparing these results to the L\'evy walk model provides a quantitative benchmark. The ergodicity breaking in the point-vortex system is of a similar magnitude to that of the L\'evy walks. For instance, the $N=40$ system ($EB \approx 0.45$) exhibits ergodicity breaking comparable to a L\'evy walk with a tail exponent of $\beta=1.5$ ($EB \approx 0.48$). This suggests that while the underlying memory mechanisms differ, the degree of transport heterogeneity in the dense chaotic flow is consistent with that produced by a moderately superdiffusive, memoryless stochastic process.

\subsection{Mean squared displacement and transport regimes}

The primary characterization of transport was performed by analyzing the ensemble-averaged mean squared displacement (EAMSD). Fitting a power law, $\langle \mathbf{r}^2(\Delta t) \rangle \sim (\Delta t)^\alpha$, to the EAMSD curves reveals a striking non-monotonic relationship between the anomalous diffusion exponent $\alpha$ and the vortex density $N$.

For sparse to intermediate densities, the transport is near-ballistic. We find $\alpha = 2.008 \pm 0.000$ for $N=5$, $\alpha = 2.032 \pm 0.001$ for $N=10$, and $\alpha = 2.100 \pm 0.001$ for $N=20$. The values of $\alpha \approx 2$ indicate that tracers are advected coherently over long distances by the large-scale velocity field, with little diffusive spreading.

However, a distinct transition occurs in the densest configuration. For $N=40$, the exponent drops significantly to $\alpha = 1.680 \pm 0.014$, signaling a clear shift from near-ballistic motion to a superdiffusive regime. This non-monotonic behavior provides a key insight into the system's dynamics. At low to intermediate $N$, the flow structure is dominated by coherent vortices that lead to prolonged, straight-line flights. At high $N$, the increased vortex-vortex interactions create a more complex and rapidly evolving chaotic sea. This environment enhances tracer trapping and accelerates velocity decorrelation, as seen in the residence time and VACF analyses. The combination of these effects disrupts the long ballistic flights, resulting in a transport process that is still superdiffusive ($\alpha > 1$) but significantly less efficient than ballistic motion ($\alpha < 2$).

\end{document}
                